% !PRINTED 186
\weberpage{186}
For if we set

\[
D = \begin{vmatrix} a_1 - 1, & a_2, & a_3\\ b_1, & b_2 - 1, & b_3\\ c_1, & c_2, & c_3 - 1 \end{vmatrix},
\]

then from the relations (4) we obtain:

\[
|A|\, D = - D,
\]

and thus, if (5) holds, $D = 0$. Accordingly, three quantities
$\lambda, \mu, \nu$ can be determined from the equations:

\weq{6}{\begin{array}{l}
\lambda = a_1 \lambda + a_2 \mu + a_3 \nu\\
\mu = b_1 \lambda + b_2 \mu + b_3 \nu\\
\nu = c_1 \lambda + c_2 \mu + c_3 \nu\\
\lambda^2 + \mu^2 + \nu^2 = 1
\end{array}}

and these three quantities $\lambda, \mu, \nu$ determine the direction of a straight line which makes the same
angles with the axes $y_1, y_2, y_3$ as with the axes $x_1, x_2, x_3$. A line laid in this
direction through the origin of coordinates is
the axis of rotation.

In the case $|A| = - 1$ this conclusion no longer applies.

Besides this there is yet another interpretation of the orthogonal
substitution $(y) = A (x)$, according to which $(x)$ and $(y)$ are the coordinates
of two different points $x$ and $y$, referred to one and
the same coordinate system. When $x$ sweeps over a part of space (a line,
surface or body), the corresponding
point $y$ fills a congruent part of space (if $|A| = - 1$, one
that is its mirror image). This is shown by the following consideration.

The group of the \gesperrt{proper} orthogonal substitutions
is equivalent to a group that can be formed from the different
positions of a body rotatable about a fixed point.
This group is obtained by taking an arbitrary position $E$ as the unit, from which one passes to any other
position $A$ by a rotation about a definite axis through
a definite angle. If $B$ is a third position, then by
the composed position $AB$ is to be understood the position
one obtains by carrying out the rotation that led to $A$
not from the unit position, but from the
position $B$. For if one takes a system of axes standing in rigid
connection with the body, e.g. the system
of the principal axes of inertia, and denotes by $(x)$ the coordinates
%CONT>>
